기룡이는 밤늦게 게임 개발을 하다가 잠이 들고 말았다. 
눈을 뜬 기룡이는 자신이 건물에 혼자 남아 있다는 사실을 깨달았다.
하지만 기룡이는 금방 무언가 이상함을 눈치챌 수 있었다.
벽들은 미로처럼 뒤엉켜 있었고, 문 너머에는 같은 구조의 방이 끝없이 반복되고 있었던 것이다.

기룡이는 자신이 만들던 게임 `루프 더 룸' 속 캐릭터가 되었음을 직감했다. 
몸은 격자를 따라 끊기듯 상하좌우로 한 칸씩 이동하였으며, 벽이 있는 칸으로는 이동할 수 없었다.

또한 무한한 루프답게 각 방의 경계를 넘어 이동하면 방의 반대편 끝 칸으로 이어진다. 
경계를 넘는 이동 규칙은 다음과 같다.
\begin{itemize}
  \item 좌우 이동: 같은 행의 가장 왼쪽 칸과 가장 오른쪽 칸이 모두 벽이 아니라면, 가장 왼쪽 칸에서 왼쪽으로 이동하여 가장 오른쪽 칸으로 갈 수 있다. 반대로 가장 오른쪽 칸에서 오른쪽으로 이동하여 가장 왼쪽 칸으로 갈 수도 있다.
  \item 상하 이동: 같은 열의 가장 위쪽 칸과 가장 아래쪽 칸이 모두 벽이 아니라면, 가장 위쪽 칸에서 위로 이동하여 가장 아래쪽 칸으로 갈 수 있다. 반대로 가장 아래쪽 칸에서 아래로 이동하여 가장 위쪽 칸으로 갈 수도 있다.
\end{itemize}

불행 중 다행으로 기룡이는 천장과 바닥까지 루프로 만들기 전에 잠들어 버렸기에 내려가는 계단만 찾으면 이곳에서 빠져나갈 수 있었다.

한 번 이동할 때마다 $1$초의 시간이 걸린다 할 때, 기룡이가 시작 위치에서부터 계단까지 가는 데 걸리는 최소 시간을 구하자.